\section{Review Questions, Scope, and Evidence Methodology}
\label{sec:methods}

\subsection{Review objective and research questions}

This study treats optimization research as an evidence-organization problem rather than as a catalogue of algorithm names. The review asks which algorithm families have been studied, which computational mechanisms distinguish them, what theoretical and empirical evidence supports their use, how benchmark design conditions the conclusions that can be drawn, which implementations are reproducible, and where the available evidence remains incomplete. Particular attention is given to structural overlap among nominally different methods, condition-dependent performance deterioration, and the distinction between documented limitations and genuinely unresolved research gaps.

Nine review questions organize the analysis. They concern algorithm and problem-class coverage; mechanism-level distinctions; evidence on quality, convergence, robustness, scalability, computational cost, and reliability; benchmark adequacy; structural similarity; condition-dependent failure regimes; reproducibility; evidence-supported gaps; and the circumstances under which a gap would justify an established remedy, a transferable mechanism, or a new optimizer. A new optimizer is therefore not a predetermined outcome of the review.

\subsection{Evidence classes and claim discipline}

Four evidence classes are kept separate throughout the study: documented facts supported directly by a source; empirical results reported by prior studies; cross-study synthesis produced by the review; and project-generated computational findings. Literature-reported evidence is not relabelled as project validation, and project results are not used to strengthen claims beyond the conditions actually evaluated.

The review uses a 90-entry algorithm census as an identity and coverage layer. The census is not interpreted as evidence for 90 distinct mechanisms. Algorithms are instead normalized into a mechanism-first representation. Following the final evidence audit, the representative mechanism layer contains 42 source-backed fingerprints linked to 46 verified source records. Named variants are retained when they introduce a relevant structural distinction; nominally different methods are not promoted to separate mechanisms merely because they use different metaphors or terminology.

\subsection{Search, screening, and evidence extraction}

The evidence-mining protocol covers exact and mathematical optimization, gradient and proximal methods, derivative-free search, evolutionary and population-based optimization, swarm methods, probabilistic and distribution-based search, Bayesian and surrogate optimization, multi- and many-objective optimization, constrained and robust optimization, dynamic and online optimization, mixed-variable and combinatorial methods, bilevel optimization, and learning-assisted optimization.

Search concepts combine optimization terminology with mechanism, benchmark, scalability, robustness, constraint, noise, dynamics, mixed-variable, reproducibility, structural-bias, comparison, sensitivity, ablation, and parameter-control terms. Sources are eligible when they contribute an original algorithm or substantial mechanism, a theoretical result, benchmark methodology, a systematic or critical synthesis, comparative evidence, reproducibility evidence, structural-bias analysis, benchmark-suite design, negative evidence, or a real-world evaluation. Sources that do not provide an inspectable algorithmic definition, cannot be linked to a verifiable bibliographic record, report selected best runs without adequate context, or use comparisons whose budgets cannot be interpreted fairly are excluded or isolated rather than silently incorporated.

For each representative algorithm, evidence extraction records identity and provenance, representation, initialization, information sources, update geometry, selection or replacement, memory, adaptation, diversity or restart mechanisms, constraint handling, derivative information, surrogate use, stopping conditions, theoretical guarantees and assumptions, benchmark coverage, parameter policy, statistical practice, implementation availability, and known limitations. Untested conditions are recorded as evidence gaps; absence of evidence is not interpreted as algorithmic failure.

\subsection{Mechanism-first normalization}

Each representative method is encoded using a common structural fingerprint comprising state type, information source, proposal geometry, memory, adaptation, selection, diversity control, constraint handling, variable scope, objective scope, derivative order, model type, parallelism, dominant overhead, and a documented or hypothesized failure signal. This normalization provides a common language for comparing methods that originate in different research traditions.

Structural similarity is interpreted as similarity of encoded computational mechanisms, not as algorithmic identity, plagiarism, or equivalent empirical performance. Similarity analyses are therefore used as an exploratory synthesis tool and are checked against known genealogical relationships and manual mechanism inspection.

\subsection{Benchmark and comparison protocol}

The computational component uses the noiseless COCO BBOB suite and retains all 24 functions rather than selecting functions after observing results. The core design considers dimensions \(D\in\{2,5,10,20,40\}\). The principal cross-dimensional campaigns use dimensions 5, 10, 20, and 40. Discovery uses BBOB instances 1--5, while a separately frozen held-out characterization uses instances 6--10. Stochastic configurations use seeds 11, 23, and 37. Evaluation budgets are normalized as \(100D\), \(300D\), and \(1000D\) objective evaluations.

The maintained comparator panel comprises Random Search as a transparent reference baseline, SciPy differential evolution with an explicitly frozen configuration and polishing disabled, bounded SciPy Nelder--Mead as a direct-search comparator, and pycma CMA-ES through an ask--tell interface. Algorithms whose publication implementation or configuration provenance did not pass the same acceptance process were not inserted merely to enlarge the comparison panel.

All publication computations use objective-level accounting. The objective recorder increments once for each actual objective call, maintains best-so-far values, records declared checkpoints, and is cross-checked against COCO's evaluation counter. Evaluation budgets are therefore defined in objective calls rather than generations or optimizer-specific iterations. Failures and unobserved checkpoints are retained rather than imputed as successful observations.

\subsection{Statistical analysis and interpretation}

The atomic experimental unit is algorithm/configuration \(\times\) benchmark function \(\times\) instance \(\times\) dimension \(\times\) seed \(\times\) budget. Primary views include target attainment, evaluations to target, best-so-far values at declared budgets, anytime or empirical cumulative distribution views where supported by suite semantics, and computational overhead as a separate cost dimension. Aggregate ranks, when reported, are descriptive and are not used to establish universal superiority.

Comparisons are matched on function, instance, dimension, seed, and budget whenever the algorithms permit such pairing. Distributional summaries and confidence intervals accompany the comparisons. Paired effect estimates with bootstrap confidence intervals are used where appropriate, and multiplicity-controlled post-hoc comparisons follow an omnibus analysis when inferential comparisons are warranted. Holm-type family-wise error control is the default for planned post-hoc families unless the endpoint structure requires a separately declared procedure.

\subsection{Discovery, held-out characterization, and frontier claims}

Discovery and confirmation data are separated to prevent a condition identified from the data from serving as its own sole confirmation. The EF003 discovery campaign uses instances 1--5. EF005 repeats the frozen algorithms, dimensions, seeds, and evaluation budgets on held-out instances 6--10. The computational and provenance pipelines, including algorithm-separated COCO logging and joint \texttt{cocopp} postprocessing, passed the frozen validation workflow.

The original frontier protocol required an exact target subset and a numerical reliability threshold to be fixed before held-out outcomes were inspected. Those exact numerical choices were not frozen before the held-out campaign was executed. Consequently, the paper does not promote a threshold-dependent binary failure frontier after the fact. The held-out campaign is used for condition-level target-runtime and scaling characterization and for assessing whether qualitative patterns persist on disjoint instances. No post hoc exact boundary, universal failure claim, or mechanism-repair claim is inferred solely from performance deterioration.

\subsection{Gap validation and remedy rule}

Evidence gaps are classified as theoretical, mechanistic, structural-bias, benchmark, scalability, robustness, constraint, dynamic, mixed-variable, reproducibility, comparison, real-world, statistical, or computational-cost gaps. A gap record distinguishes what is documented, partially supported, contradictory, untested, or project-validated.

A new mechanism or optimizer is considered only when evidence supports a mechanism-level repair opportunity that cannot be adequately explained by tuning, budget normalization, an existing adaptive variant, restart, established constraint handling, or another known remedy. Because the frozen evidence does not establish such a confirmatory G4/G5 repair opportunity, no new optimizer or result-driven ablation campaign is introduced in this paper.

\subsection{Reproducibility and evidence freeze}

The research workflow separates calibration, pilot, rehearsal, run-complete-unverified, and project-validated states. Calibration, pilot, and rehearsal outputs are never cited as publication evidence. Benchmark software, comparator configurations, target semantics, statistical rules, execution stacks, and interpretation boundaries are versioned in the repository.

Before manuscript drafting, the P031 evidence-freeze workflow validated the execution-freeze records, literature/source integrity, final pre-freeze audit, and a Git-blob manifest of the principal evidence and publication-control artifacts. The manuscript is written against this frozen boundary. Subsequent scientific changes require a separately versioned amendment and cannot silently replace the evidence supporting the reported conclusions.
